\documentclass[aps,prl,twocolumn,showpacs,superscriptaddress,groupedaddress]{revtex4-2}
\usepackage{graphicx,amsmath,amssymb,bm}

\begin{document}

\title{Vibronic-Mediated Decoherence: A Microscopic Foundation for Temperature-Dependent Dephasing in DTQEM}

\author{Reddouane Berramdane}
\affiliation{Independent Researcher, Oujda, Morocco}
\email{reddoma@gmail.com}

\date{June 2026}

\begin{abstract}
We present a microscopic foundation for the path-dephasing coefficient in the DTQEM framework by deriving its temperature dependence from a shared vibronic mode. Starting from a minimal pure-dephasing Hamiltonian and applying the exact independent-boson solution, we obtain an effective dephasing coefficient of the form
\[
a_{\mathrm{eff}}(T)=a_0+2S\coth\!\left(\frac{\hbar\omega_v}{2k_B T}\right),
\]
where \(S=g^2/\omega_v^2\) is the Huang–Rhys factor and \(\omega_v\) is the dominant vibrational frequency. This expression replaces the phenomenological bilinear cross-coupling of earlier DTQEM versions (Papers I and II), recovers the high-temperature linear behavior, predicts a low-temperature saturation, and yields a falsifiable isotope shift through \(\omega_v\propto \mu^{-1/2}\). We propose tests in molecular interferometry, NV centers in diamond, and other open quantum systems with a dominant vibronic mode.
\end{abstract}

\maketitle

\section{Introduction}

Open quantum systems often exhibit decoherence that depends not only on environmental noise but also on the internal vibrational structure of the system. In path-based interference settings, this suggests that the effective dephasing strength may acquire a temperature dependence governed by vibronic coupling. In earlier DTQEM formulations \cite{PaperI,PaperII}, the joint effect of path information and temperature was captured by a bilinear term \(c\,I_{\mathrm{path}}\Delta T/T_{\mathrm{ref}}\). However, that term was treated as a free parameter without an explicit connection to a molecular or solid-state vibrational mode. Here we show that this bilinear term is the first-order local expansion of a temperature-dependent dephasing coefficient \(a(T)\). We replace the phenomenological ansatz with a first-principles derivation based on the independent-boson model, which is exactly solvable in the pure-dephasing limit. The resulting theory provides a direct link between the effective coefficient \(a(T)\), the Huang–Rhys factor \(S\), and the dominant vibrational frequency \(\omega_v\).

\section{Microscopic model}

We consider a two-level system coupled to a single dominant vibrational mode through pure dephasing,
\[
H=\frac{\Delta}{2}\sigma_z+\hbar\omega_v b^\dagger b+\hbar g\,\sigma_z(b+b^\dagger),
\]
which is the standard independent-boson model. This minimal Hamiltonian captures the essential vibronic contribution to dephasing when the cross-spectral density is concentrated near a single frequency \(\omega_v\). For an initially factorized state with the bosonic mode in thermal equilibrium, the off-diagonal coherence is obtained exactly as
\[
C(t)=C_0\exp\!\left[-2S\coth\!\left(\frac{\hbar\omega_v}{2k_B T}\right)(1-\cos\omega_v t)\right],
\]
where
\[
S=\frac{g^2}{\omega_v^2}
\]
is the Huang–Rhys factor. The oscillatory term \((1-\cos\omega_v t)\) produces coherence revivals at \(t=2\pi m/\omega_v\). The temperature-dependent effective dephasing coefficient is derived microscopically in Appendix A, where we show that the exact single-mode dynamics reduces, after coarse-graining over the measurement window, to
\[
a_{\mathrm{eff}}(T)=a_0+2S\coth\!\left(\frac{\hbar\omega_v}{2k_B T}\right).
\]
Here \(a_0\) denotes the residual dephasing floor at zero temperature, representing temperature-independent dephasing mechanisms not captured by the single-mode vibronic contribution, such as measurement backaction, static disorder, or other bath modes. In the minimal independent-boson model, \(a_0=0\); a non-zero \(a_0\) accounts for additional processes present in a realistic experiment. In the low-temperature limit \(k_B T\ll \hbar\omega_v\), the coefficient saturates to \(a_{\mathrm{eff}}(T)\to a_0+2S\), while in the high-temperature limit \(k_B T\gg \hbar\omega_v\), it becomes \(a_{\mathrm{eff}}(T)\approx a_0+4Sk_B T/(\hbar\omega_v)\). These limiting behaviors recover the expected saturation and linear thermal growth, respectively.

\section{Predictions and limiting behaviors}

The microscopic form of the effective dephasing coefficient yields three immediate predictions. First, the low-temperature behavior is saturating:
\[
a_{\mathrm{eff}}(T)\to a_0+2S \qquad (k_B T\ll \hbar\omega_v).
\]
Second, the high-temperature regime is linear in temperature:
\[
a_{\mathrm{eff}}(T)\approx a_0+\frac{4Sk_B}{\hbar\omega_v}T \qquad (k_B T\gg \hbar\omega_v).
\]
Third, the crossover temperature scales with the dominant vibrational frequency as
\[
T^*\sim \frac{\hbar\omega_v}{2k_B}.
\]
Since \(\omega_v\propto \mu^{-1/2}\) for a vibrational mode with reduced mass \(\mu\), isotope substitution leads to a measurable shift in the crossover scale:
\[
\frac{T^{*\prime}}{T^*}=\sqrt{\frac{\mu}{\mu'}}.
\]
Equivalently, the effective dephasing strength changes as
\[
a'_{\mathrm{eff}}(T)=a_0+2S'\coth\!\left(\frac{\hbar\omega'_v}{2k_B T}\right),\qquad S'=\frac{g^2}{\omega_v'^2}.
\]
This isotope dependence is the most distinctive experimental signature of the model, because it ties the temperature dependence of dephasing directly to a measurable vibrational frequency. In systems with a single dominant mode, the predicted shift can be extracted from temperature-dependent interference or coherence-decay measurements. In multimode environments, the same logic generalizes by replacing the single frequency \(\omega_v\) with an effective spectral density.

\section{Numerical estimates for C$_{70}$ fullerenes}

As an illustration, consider the C$_{70}$ molecule studied in the matter-wave interferometry experiments of Hornberger et al.~\cite{Hornberger2003}. The dominant vibrational mode contributing to vibronic dephasing is the ``squashing mode'' at approximately \(\omega_v/2\pi\approx 27\) THz (\(900\ \mathrm{cm^{-1}}\)). This gives a crossover temperature \(T^*=\hbar\omega_v/(2k_B)\approx 195\) K. An illustrative estimate based on the room-temperature visibility trend suggests an effective Huang–Rhys factor \(S\) of order \(0.8\); a more precise determination requires temperature-dependent visibility data for confirmation. For isotopically substituted \(^{13}\)C$_{70}$, or for comparison between \(^{12}$C$_{60}$ and \(^{13}$C$_{60}$, the reduced-mass change leads to a relative shift \(\Delta T^*/T^*\approx -3.9\%\), which is experimentally accessible with modern cryogenic interferometry.

\section{Comparison with experiment}

The present model can be tested in any system where a dominant vibronic mode controls the loss of coherence. A natural first target is matter-wave interferometry with large molecules, where the interference visibility is measured as a function of temperature and environmental conditions. In particular, fullerene experiments provide a useful benchmark because the coherence signal is directly accessible and the vibrational spectrum is sufficiently structured to support a reduced single-mode description in first approximation. Within the DTQEM framework, the measured visibility may be written schematically as
\[
V(T)\approx V_0\exp[-a_{\mathrm{eff}}(T)I_{\mathrm{path}}],
\]
with \(a_{\mathrm{eff}}(T)=a_0+2S\coth(\hbar\omega_v/2k_B T)\). At fixed temperature, the fit determines an effective coefficient \(a_{\mathrm{eff}}\), while temperature-dependent measurements allow the separation of \(a_0\), \(S\), and \(\omega_v\).

A second promising platform is the NV center in diamond, where temperature-dependent dephasing is routinely measured and the phonon environment is well characterized. In such systems, the low-temperature saturation and high-temperature linear growth predicted by the model can be checked directly by fitting the coherence time \(T_2\) or the Ramsey contrast as a function of temperature. The same framework also predicts an isotope-dependent shift in the crossover temperature, which could be probed by isotopically purified or enriched samples. More generally, the model is most informative when the experiment provides data across a sufficiently broad temperature range. Single-temperature measurements can constrain only \(a_{\mathrm{eff}}\), whereas multi-temperature data are needed to test the thermal functional form and to determine whether a single dominant mode is indeed responsible for the observed dephasing. If the data require multiple frequencies, the model should be generalized to a spectral-density formulation.

\section{Discussion}

The main conceptual advance of the present formulation is that it replaces a purely phenomenological temperature dependence by a microscopic quantity with a direct physical interpretation. The parameter \(S\) is measurable independently from optical spectroscopy, which makes the theory more predictive than a free bilinear coefficient. The model retains a small number of parameters \((a_0,S,\omega_v)\) and reproduces the earlier bilinear DTQEM form as a local first-order approximation around a reference temperature.

A key limitation of the present model is the single-mode approximation. For systems with a broad spectral density, the effective coefficient should be replaced by
\[
a_{\mathrm{eff}}(T)=a_0+\int_0^\infty J_{\mathrm{HR}}(\omega)\coth\!\left(\frac{\hbar\omega}{2k_B T}\right)d\omega,
\]
where \(J_{\mathrm{HR}}(\omega)\) is the Huang–Rhys spectral density, with the normalization absorbed into the definition of \(J_{\mathrm{HR}}\). This general expression reduces to the single-mode form when \(J_{\mathrm{HR}}(\omega)=2S\delta(\omega-\omega_v)\). In multimode cases, the single-mode form may still be useful for identifying a dominant feature, but the more general spectral integral must be retained for quantitative accuracy.

\section{Conclusion}

We have derived a temperature-dependent effective dephasing coefficient from the exact independent-boson solution of a single-mode pure-dephasing model. The resulting form,
\[
a_{\mathrm{eff}}(T)=a_0+2S\coth\!\left(\frac{\hbar\omega_v}{2k_B T}\right),
\]
provides a microscopic basis for DTQEM, predicts low-temperature saturation, high-temperature linear growth, and an isotope-dependent crossover temperature. This makes the model directly testable in molecular interferometry, solid-state qubits, and related open-quantum-system platforms.

\begin{acknowledgments}
The author thanks the open-source communities behind NumPy, SciPy, and Matplotlib. AI language models (DeepSeek, Claude, Qwen) were used as analytical and editorial tools under full human scientific supervision. All scientific claims are the sole responsibility of the author.
\end{acknowledgments}

\appendix
\section{Coarse-graining to an effective dephasing coefficient}

The exact off-diagonal coherence for a two-level system coupled to a single bosonic mode is
\[
C(t)=C_0\exp\!\left[-2S\coth\!\left(\frac{\hbar\omega_v}{2k_B T}\right)(1-\cos\omega_v t)\right],
\]
where
\[
S=\frac{g^2}{\omega_v^2}
\]
is the Huang–Rhys factor. The factor \((1-\cos\omega_v t)\) oscillates between 0 and 2 with period \(2\pi/\omega_v\), producing periodic coherence revivals at times \(t=2\pi m/\omega_v\), with \(m\in\mathbb{Z}\).

In an actual measurement, the observation time \(\tau_{\mathrm{meas}}\) typically spans many vibrational periods, so that \(\omega_v \tau_{\mathrm{meas}}\gg 1\). The detector therefore probes a coarse-grained coherence envelope rather than the instantaneous oscillatory signal. The effective coefficient \(a_{\mathrm{eff}}(T)\) is not obtained by a simple time average of the instantaneous rate; rather, it is extracted by fitting the coarse-grained coherence envelope (or the cycle-averaged visibility) to a single exponential decay. This definition is operational and corresponds to a fit to finite-resolution experimental data.

In the weak-coupling regime, \(2S\coth(\hbar\omega_v/2k_B T)\ll 1\), the observed fringe visibility may be approximated by a monotonic exponential with an effective coefficient
\[
a_{\mathrm{eff}}(T)=a_0+2S\coth\!\left(\frac{\hbar\omega_v}{2k_B T}\right),
\]
where \(a_0\) accounts for temperature-independent dephasing channels not captured by the single-mode vibronic contribution.

More generally, the exact cycle average gives a modified-Bessel correction:
\[
\frac{\langle C\rangle}{C_0}=e^{-A}I_0(A),\qquad A=2S\coth\!\left(\frac{\hbar\omega_v}{2k_B T}\right),
\]
where \(I_0\) is the modified Bessel function of the first kind. This correction reduces to \(\langle C\rangle/C_0\approx e^{-A}\) in the weak-coupling limit \(A\ll 1\), but becomes important when \(S\gtrsim 1\).

The effective form \(a_{\mathrm{eff}}(T)\) is valid when the measurement does not resolve individual coherence revivals. It breaks down for ultrafast interrogation, strong coupling, or multimode environments, in which case the full time-dependent expression must be retained. In the high-temperature limit \(k_B T\gg \hbar\omega_v\), one recovers \(a_{\mathrm{eff}}(T)\approx a_0+4Sk_B T/(\hbar\omega_v)\), while in the low-temperature limit \(k_B T\ll \hbar\omega_v\), the coefficient saturates to \(a_{\mathrm{eff}}(T)\to a_0+2S\).

\begin{thebibliography}{99}

\bibitem{Mahan}
G. D. Mahan,
\emph{Many-Particle Physics}, 3rd ed.
(Plenum, New York, 2000).

\bibitem{Leggett}
A. J. Leggett \emph{et al.},
Rev. Mod. Phys. \textbf{59}, 1 (1987).

\bibitem{Breuer}
H.-P. Breuer and F. Petruccione,
\emph{The Theory of Open Quantum Systems}
(Oxford University Press, 2002).

\bibitem{Weiss}
U. Weiss,
\emph{Quantum Dissipative Systems}, 4th ed.
(World Scientific, 2012).

\bibitem{Hornberger2003}
K. Hornberger \emph{et al.},
Phys. Rev. Lett. \textbf{90}, 160401 (2003).

\bibitem{PaperI}
R. Berramdane,
\textit{DTQEM: A Phenomenological Model for Correlated Path-Information and Thermal Decoherence (v18.0-C)},
Zenodo (2026), DOI:10.5281/zenodo.20670032.

\bibitem{PaperII}
R. Berramdane,
\textit{Microscopic Derivation of Joint-Bath Decoherence: Cross-Spectral Density and Inter-Bath Correlation},
Zenodo (2026), DOI:10.5281/zenodo.20670032 (companion paper).

\end{thebibliography}

\end{document}
